\scanpage{10}
\uline{N.B.} Cela signifie donc que \deleted{les composantes} \inserted{l'ens des} composantes irréductibles relatives de $f^{-1}(y')$ \deleted{est \unsure{contenu}} s'\uline{injecte} dans celui des composantes irréductibles \deleted{\unsure{absolues|distinctes}} \inserted{relatives} de $f^{-1}(\eta')$\inserted{\unsure{par contre|par division|qui est}}, ce n'est pas néc. surjectif. On a cependant ceci :

\uline{Cor} Pour que $M$ soit un \uline{\unsure{revêtem\textsuperscript{t}|revêt\textsuperscript{mt}|ouvert}} étale de $Y$, il f et suffit que l'on ait de plus, avec les notations précédentes :

(iii) \deleted{\unsure{Toute|Pour}} Pour \deleted{toute} \inserted{p\textsuperscript{t} gén. $\xi'$ d'une} composante irréductible relative \deleted{\unsure{$\xi$ de $\xi'$}} de $ {f'}^{-1}(\eta')$, il existe un p\textsuperscript{t} gén. $x'$ d'une comp. irréd. relative de $ {f'}^{-1}(y')$ tel que $x'\in\overline{\xi'}$.

[Cette condition n'est pas tjrs vérifiée sous la seule condition que $f:X\to Y$ soit propre, car \deleted{\unsure{si une composante}} les fibres de $\overline{\xi'}$ peuvent être de dim str. inférieure à celle des composantes de $ {f'}^{-1}(y')$ qui les contiennent. On verra que ceci n'est pas possible s'il n'y a pas de \unsure{'idéaux|idéaux} premiers immergés ---]
% archive sheet number: 9
